\documentclass[11pt]{article}

\usepackage[margin=1in]{geometry}
\usepackage{amsmath,amssymb}
\usepackage[T1]{fontenc}

\title{Random Sparse Covariance Generation}
\author{}
\date{}

\begin{document}

\maketitle

Let $p$ be the number of features and let $d$ be the target average number of
nonzero off-diagonal entries per row. We sample
\[
M=\operatorname{round}(pd/2)
\]
distinct unordered feature pairs uniformly at random. For each selected pair,
we independently sample a positive weight from $\operatorname{Uniform}(0.5,1.5)$.
Placing these weights symmetrically produces a sparse symmetric matrix $A$ with
zero diagonal. Thus, the average number of nonzero off-diagonal entries per row
is approximately $d$, without imposing a block structure.

For a fixed $0<\rho<1$, define
\[
\Sigma=I_p+\rho\frac{A}{\lVert A\rVert_{\mathrm{op}}}.
\]
Because every eigenvalue of $A/\lVert A\rVert_{\mathrm{op}}$ lies in
$[-1,1]$,
\[
\lambda_{\min}(\Sigma)\geq 1-\rho>0.
\]
Therefore, every generated matrix is a valid covariance matrix. We accept a
draw only when
\[
\lambda_1(\Sigma)-\lambda_2(\Sigma)\geq \Delta,
\]
where $\Delta>0$ is specified before the search. If no draw satisfies this
condition within the maximum number of attempts, the procedure stops with an
error.

All generation, validation, and search functions are contained in
\texttt{src/001\_find\_gap\_covariance.R}. Its main interface takes $p$, $d$,
$\rho$, $\Delta$, the maximum number of attempts, and a random seed. It returns
the accepted covariance matrix, its leading eigenpairs, the generation
settings, and sparsity and eigenvalue diagnostics. The script
\texttt{code/001\_find\_sparse\_covariance\_matrix.R} declares its input and
output filenames at the beginning, calls this interface for $p=100$ and
$p=1000$, and saves the results as
\texttt{data/raw/random\_sparse\_covariance\_p100.rds} and
\texttt{data/raw/random\_sparse\_covariance\_p1000.rds}. The computational
module does not perform file input or output. We use $d=8$ and $\rho=0.9$ for
both $p=100$ and $p=1000$. Relative to a denser graph with the same bounded
spectral perturbation, this choice produces fewer but stronger nonzero
covariances and is therefore more suitable for evaluating sparse covariance
estimation.

\end{document}
